\begin{exo}[Un exercice classique]
	\label{seriesfonc1}
	Soit $(f_n)_n$ une suite de fonctions continues convergeant uniformément 
	sur $\R$ vers $f$. 
	Soit $(x_n)_{n\in\N}$ une suite de réels convergeant vers $x$.\\
	Montrer que la suite $(f_n(x_n))_{n\in\N}$ converge et calculer sa limite.
\end{exo}

\begin{solution}
	$f$ étant limite uniforme d'une suite de fonctions continues, elle est continue. 
    Il vient alors que \[f(x_n)\xrightarrow[n\to+\infty]{} f(x)\tag{1}\]
	Puis, pour $n\in\N$, on a 
    \begin{align*}
        |f_n(x_n)-f(x)|&=|f_n(x_n)-f(x_n)+f(x_n)-f(x)|\\ 
                    &\leq|f_n(x_n)-f(x_n)|+|f(x_n)-f(x)|
    \end{align*}
	Or, puisqu'il y a convergence uniforme,
    \[
        |f_n(x_n)-f(x_n)|\leq \norm{f_n-f}_\infty^{\R}\xrightarrow[n\to+\infty]{} 0
    \] 
    Puis (1) donne que $|f(x_n)-f(x)|\xrightarrow[n\to+\infty]{}0$.
    Par encadrement, on a $|f_n(x_n)-f(x)|\xrightarrow[n\to+\infty]{} 0$, soit \[\boxed{f_n(x_n)\xrightarrow[n\to+\infty]{}f(x)}\]
\end{solution}

\begin{exo}[L'exponentielle comme limite uniforme de polynômes ?]
	\label{seriesfonc2}
	Existe-t-il une suite de polynômes convergeant uniformément sur $\R$ vers $\exp$ ?
\end{exo}

\begin{solution}
	Soit $(p_n)_{n\in\N}$ une suite de polynômes convergeant uniformément vers $\exp$ sur $\R$, soit 
    \[\forall \varepsilon > 0,\ \exists n_0\geq 0,\ \forall n\geq n_0,\ \forall x\in\R,\ |p_n(x)-\exp(x)|\leq \varepsilon\]
	Ainsi, si on "applique" cette phrase pour $\varepsilon = 333$, on a l'existence de $n_0\geq 0$ tel que, pour $n\geq n_0$ et pour $x\in\R$, on ait $|p_n(x)-\exp(x)|\leq 333$. 
    Mézalor \[|p_n(x)-p_{n_0}(x)|\leq |p_n(x)-\exp(x)|+|p_{n_0}(x)-\exp(x)|\leq 666\]
	Ceci valant pour tout $n\geq n_0$ et pour tout $x\in\R$, on a que, pour tout $n\geq n_0$, le polynôme $p_n-p_{n_0}$ est borné sur $\R$, donc constant, 
    d'où l'existence d'une suite de réels $(\alpha_n)$ telle que \[\forall n\geq n_0,\ \forall x\in\R,\ p_n(x)-p_{n_0}(x)=\alpha_n\]
	La suite $(p_n(42))_{n\in\N}$ convergeant, puisque CU implique CS, on a que la suite $(\alpha_n)$ converge également, 
    il suffit d'évaluer la précédente expression en $42$ ; on note $\alpha$ sa limite. Pour $x\in\R$, on a \[\forall n\geq n_0,\ p_n(x)-p_{n_0}(x)=\alpha_n\]
	On passe à la limite quand $n$ tend vers $+\infty$ et on obtient \[\exp(x)=p_{n_0}(x)+\alpha\]
	autrement dit, $\exp$ est un polynôme, ce qui est exclu.
\end{solution}

\begin{exo}[Étude de continuité]
    \label{seriesfonc3}
    On considère la fonction définie par 
    \[f(x):=\sum_{n=1}^{+\infty}\frac{n^x}{x^n}\]
    Déterminer le domaine $D$ de définition de $f$ et étudier la continuité de $f$ sur $D$.
\end{exo}

\begin{solution}
    Pour $|x|>1$, on voit facilement que $|n^x/x^n|=o(1/n^2)$, donc par comparaison de séries à termes positifs, la série converge absolument.
    Pour $x=-1$, le terme général est $(-1)^n/n$, la convergence découle du CSSA (série harmonique alternée).
    Pour $x\in(-1,1]$, le terme général ne tend pas vers $0$ (si $x \in (0,1]$, $n^x/x^n \to +\infty$ ; si $x \in (-1,0]$, $|n^x/x^n| = n^x/|x|^n \to +\infty$). La série diverge grossièrement.
    On a alors $D=(-\infty,-1]\cup(1,+\infty)$.
    
    Pour la continuité, on étudie la convergence uniforme sur ces deux composantes connexes :
    \begin{itemize}
        \item Sur tout compact $[a,b] \subset (1,+\infty)$, on a pour tout $x \in [a,b]$, $|n^x/x^n| \leq n^b/a^n$. Comme $a>1$, $n^b/a^n = o(1/n^2)$, d'où la convergence normale (et donc uniforme) de la série de fonctions.
        \item Sur $(-\infty,-1]$, on pose $u_n(x) = n^x/x^n = (-1)^n \frac{n^x}{|x|^n}$. Pour tout $x \in (-\infty,-1]$, la suite $(|u_n(x)|)_n$ décroît vers $0$. D'après le critère spécial des séries alternées
        \[
            \forall x \in (-\infty,-1], \quad |R_N(x)| \leq |u_{N+1}(x)| = \frac{(N+1)^x}{|x|^{N+1}} \leq \frac{1}{N+1}
        \]
        Cette majoration, indépendante de $x$, prouve la convergence uniforme de la série sur \textit{tout} l'intervalle $(-\infty,-1]$.
    \end{itemize}
    La série de fonctions définissant $f$ converge uniformément sur $(-\infty, -1]$ et sur tout compact de $(1,+\infty)$. Les fonctions $x \mapsto n^x/x^n$ y étant continues, on en déduit par le théorème de continuité des séries de fonctions que $f$ est continue sur $D$.
\end{solution}